\section{An effective triple-genus reduction by exact ramification}
\label{sec:projected-triple-ramification}

\paragraph{Status and scope.}
HD-GC01--02 are new root derivations, written 22 November 2026.
Complete author and separate independent whole-proof readings pass. Root has read the entire independent audit report; its distinct orbit enumeration is recorded with its exact scope. All covers below are actual
connected smooth complex covers, with their indicated quotient maps.
No numerical constant from an asymptotic theorem is used. This result
does not assert a numerical bound in the remaining low-genus case.

\paragraph{Notation.}
Let an actual Galois cover of $\mathbf P^1$ have group $A_n$ or $S_n$,
where $n\geq8$. Let $Y_t$ be its quotient for the action on ordered
$t$-tuples of distinct letters, and $X_t$ its quotient for the action on
unordered $t$-sets, for $t=1,2,3$. Thus $Y_1=X_1$ and the natural maps
$\pi_t:Y_t\longrightarrow X_t$ have degree $t!$. All these covers are
connected. Set
\[
 g=g(Y_1),\quad h=g(Y_2),\quad b=g(X_2),\quad c=g(X_3).
\]
For a finite map $f$ of smooth curves, $R_f$ denotes its total
ramification degree $\sum_Q(e_f(Q)-1)$. For an inertia permutation $x$,
write its individual cycle lengths as $r_1,\ldots,r_k$, including
fixed-point cycles, and put
\[
 I(x)=n-k,\qquad
 A(x)=\sum_{i,j=1}^k\bigl(r_i-\gcd(r_i,r_j)\bigr).
\]
Here the index $i$ distinguishes equal-length cycles. For a positive
integer $r$, let $v_2(r)$ be its exponent of $2$.

\paragraph{HD-GC01 (exact local triple ramification).}
The contribution to $R_{Y_2/Y_1}$ above the branch point with inertia
$x$ is $A(x)$. The contribution to $R_{\pi_3}$ is
\[
 T(x)=3\!\sum_{\substack{1\leq i,j\leq k\\v_2(r_i)>v_2(r_j)}}
           \gcd(r_i,r_j)
       +4\#\{i:3\mid r_i\}.
\]
In particular
\[
                    T(x)\leq3A(x)+2I(x).
\]
These assertions also apply to the identity permutation, with all
quantities zero.

\begin{proof}
The number of $x$-orbits on ordered distinct pairs is
$\sum_{i\ne j}\gcd(r_i,r_j)+\sum_i(r_i-1)$: between two different
cycles the orbit length is their least common multiple, and within a
single cycle the nonzero displacement specifies one orbit. The
ramification chain rule for $Y_2\to Y_1\to\mathbf P^1$ therefore gives
\[
 (n-1)k-\#\operatorname{Orb}_{Y_2}(x)
 =nk-\sum_{i,j}\gcd(r_i,r_j)=A(x).
\]
This is also the $t=3$ instance of the elementary cycle-count formula
in Neftin--Zieve, Lemma 5.4; the direct calculation here fixes the
normalization without relying on any of its asymptotic constants.

Consider an $x$-orbit of unordered triples of length $L$. The
permutation $x^L$ on its three letters generates a cyclic subgroup of
$S_3$, of order $m\in\{1,2,3\}$. Above the corresponding point of
$X_3$ there are $6/m$ points of $Y_3$, each with ramification index
$m$. This orbit therefore contributes $6-6/m$ to $R_{\pi_3}$.

If $m=2$, the triple consists of an antipodal pair in an even cycle of
length $r_i$, together with one letter in another cycle of length
$r_j$. Existence of a power that interchanges the pair and fixes the
third letter is equivalent to $v_2(r_i)>v_2(r_j)$. Conversely every
such triple has $m=2$. There are $(r_i/2)r_j$ triples of this form and
their orbit length is $\operatorname{lcm}(r_i/2,r_j)$, so their number
of orbits is
\[
 \frac{(r_i/2)r_j}{\operatorname{lcm}(r_i/2,r_j)}
 =\gcd(r_i/2,r_j)=\gcd(r_i,r_j).
\]
Each contributes $3$. This counts each such triple once: its
antipodal pair, and hence the two cycles in their indicated order,
are uniquely determined by its nontrivial stabilizer on the triple.

If $m=3$, the three letters lie in one cycle of length divisible by
$3$ and are equally spaced in that cycle. Each such cycle gives
exactly one orbit of triples, contributing $4$. These alternatives
exhaust the nontrivial cyclic subgroups of $S_3$, proving the formula.

For each pair in the first sum, $\gcd(r_i,r_j)\leq r_i/2$, whence
$\gcd(r_i,r_j)\leq r_i-\gcd(r_i,r_j)$. All other summands in $A(x)$
are nonnegative. For each cycle divisible by $3$, one has
$4\leq2(r_i-1)$. Adding these two estimates proves the inequality.
\end{proof}

\paragraph{HD-GC02 (numerical genus bounds and the residual case).}
In the notation above,
\[
 12(c-b)\geq
 2(n-8)h+(6n-10)g-12n+26.                         \tag{GC1}
\]
Consequently each of the following hypotheses implies
$g(X_3)>n^2/12$:
\[
 \begin{array}{ll}
  g(Y_1)\geq2,& n\geq16,\\
  g(Y_2)\geq n,& n\geq28.
 \end{array}                                      \tag{GC2}
\]
Thus, for $n\geq28$, the task of proving a uniform numerical
quadratic lower bound for $g(X_3)$ reduces to actual covers with
$g(Y_1)\leq1$ and $g(Y_2)<n$.

\begin{proof}
Put $A=R_{Y_2/Y_1}$ and $T=R_{\pi_3}$. Summing HD-GC01 over the
actual branch points and applying Riemann--Hurwitz to
$Y_1\to\mathbf P^1$ gives
\[
       T\leq3A+4(n+g-1),\qquad
       A=2(h-1)-2(n-1)(g-1).
\]
Riemann--Hurwitz for the three actual maps
$Y_3\to Y_2$, $Y_2\to X_2$ and $Y_3\to X_3$, of respective
degrees $n-2$, $2$ and $6$, gives exactly
\[
 12(c-b)=2(n-5)(h-1)+R_{Y_3/Y_2}+3R_{\pi_2}-T.
\]
In particular
\[
 \begin{split}
 12(c-b)&\geq2(n-5)(n-1)(g-1)
                    +(n-8)A-4(n+g-1)\\
         &=2(n-8)h+(6n-10)g-12n+26,
 \end{split}
\]
which proves (GC1). No sign condition on $h-1$ was assumed when
discarding the nonnegative ramification degrees.

If $g\geq2$ and $n\geq16$, the coefficient of $g$ in the first
expression is positive, and $A\geq0$ with $n-8\geq0$. Hence
\[
       12(c-b)\geq2n^2-16n+6>n^2.
\]
The strict inequality follows from $n^2-16n+6\geq6$ at $n=16$
and its increase thereafter. If $h\geq n$ and $n\geq28$, (GC1),
$g\geq0$ and $6n-10>0$ give
\[
       12(c-b)\geq2n^2-28n+26>n^2.
\]
Here $n^2-28n+26\geq26$ at $n=28$ and increases thereafter.
In either case $b\geq0$ proves (GC2).
\end{proof}

\paragraph{Effectivity and geometric meaning.}
In either case (GC2), any required inequality $g(X_3)>G$ follows
from the additional explicit condition $n^2\geq12G$. Likewise the
symbolic triple-genus input $g(X_3)>cn^2/12$ in HD-FT09--12 can use
$c=1$ and threshold $28$ in these cases. The full surviving
low-genus class has not been bounded here. The result is a constraint
on actual quotient covers, hence on their weight-one cohomology and
possible curve suppliers. It does not construct integral cycles or
bound a detector on a new Weil component. The existing condition-A,
condition-I, projector and spreading interfaces in the paper retain
their previously recorded exact locators and hypotheses.
